\begin{lemma}[The forward slopes of the contact envelope converge to the envelope Hamiltonian]
Let $F$ be a finite-dimensional real inner-product space --- the ambient
convention of this development --- and work in the space-time
$X=\ell^2(F\times\R)$, writing a point as $z=(x,t)$ and
$P=\{z\in X:\ t>0\}$ for the positive-time half space. Put $z_*=(0,1)$. Let
$C_0\ge1$, let $\kappa>0$ and let $v:X\to\R$. Assume:
\begin{enumerate}
\item\label{it:cesl-class} $v\in\Aclass_{m,C_0}$ in the sense of
\noderef{IsFitClass}, with $m=\dim F$;
\item\label{it:cesl-cent} $v$ is centered at $z_*$ with constant $C_0$ in the
sense of \noderef{IsCenteredAt};
\item\label{it:cesl-jensen} $v$ satisfies the horizontal Jensen bound with
constant $\kappa$ on $I=\R$ in the sense of \noderef{HasJensenBound};
\item\label{it:cesl-grow} $v$ has subquadratic spatial growth on compact time
intervals of $I=\R$ in the sense of \noderef{SubquadraticGrowth};
\item\label{it:cesl-contact} $W(z_*)=0$, where $W=\Cenv_tv-\kappa^2t/2$ is the
contact envelope of \noderef{ContactEnv} formed from $v$ and $\kappa$;
\item\label{it:cesl-Wcont} $W$ is continuous on $P$;
\item\label{it:cesl-Wgrow} $W$ has subquadratic spatial growth on compact time
intervals of $I=(0,\infty)$;
\item\label{it:cesl-Wsub} $W$ is a viscosity subsolution of the limiting
equation on $P$ in the sense of \noderef{IsLimitSubOn};
\item\label{it:cesl-sand} $v(z)-\kappa^2t/2\le W(z)\le v(z)$ for every
$z=(x,t)$ with $t>0$.
\end{enumerate}
Then there is a real number $H$ such that the envelope Hamiltonian of
\noderef{EnvHamiltonian} satisfies
\[
  H_v(z_*)=H
  \qquad\text{and}\qquad
  \lim_{\ell\downarrow0}\ \Senv_\ell^+W(z_*)=H ,
\]
the limit being over $\ell\to0$ within $(0,\infty)$ and
$\Senv_\ell^+$ the forward envelope slope of \noderef{EnvSlopeP}.
\end{lemma}
\begin{proof}
Hypotheses \ref{it:cesl-contact}--\ref{it:cesl-sand} are the four properties of
the contact envelope supplied by \noderef{ContactEnvelopeProperties} that this
argument reads, together with the contact; the construction of such a $v$ in
the application is not re-narrated. Both conclusions are returned together
because $H$ is produced in the course of the proof.

Write $t(z)$ for the time coordinate and note $t(z_*)=1$ and, more generally,
$t\bigl((y,c)\bigr)=c$: these are the only facts about the coordinates used
below. By \ref{it:cesl-cent} and \noderef{IsCenteredAt},
\begin{equation}\label{eq:cesl-v0}
  v(z_*)=0 ,
\end{equation}
and by \ref{it:cesl-class} and \noderef{IsFitClass} the function $v$ is
continuous on all of space-time and is a viscosity solution of $F[v]=0$ there.
Also, by \noderef{ContactEnv} at $t=1$, hypothesis \ref{it:cesl-contact} reads
$\Cenv_1v(z_*)-\kappa^2/2=0$, so with \eqref{eq:cesl-v0}
\begin{equation}\label{eq:cesl-bary}
  \Cenv_1v(z_*)-v(z_*)=\frac{\kappa^2}{2} .
\end{equation}

\emph{Step A: $H_v(z_*)$ is a real number.} \noderef{EnvHamiltonianSolution}
applies with $I=\R$ and $f=v$, its hypotheses being the continuity just
recorded, \ref{it:cesl-grow} and the solution property, the open set being all
of space-time. Its third clause gives $H_v(z_*)\ne+\infty$, and its second
clause, whose hypothesis $1+\sigma\in\R$ is vacuous here, gives
\begin{equation}\label{eq:cesl-ham}
  \inf_{\ell>0}\ \Senv_\ell^-v(z_*)=H_v(z_*) ,
\end{equation}
with $\Senv_\ell^-$ the backward envelope slope of \noderef{EnvSlopeM}. By
\ref{it:cesl-class}, \ref{it:cesl-cent}, \ref{it:cesl-jensen}, $\kappa>0$ and
\eqref{eq:cesl-bary}, the hypotheses of \noderef{ContactFiniteSlope} hold, so
$H_v(z_*)\ne-\infty$. A value of $\overline\R$ that is neither $+\infty$ nor
$-\infty$ is a real number; let $H$ be its real value, so that
\begin{equation}\label{eq:cesl-H}
  H_v(z_*)=H .
\end{equation}
That is the first conclusion. \noderef{ContactFiniteSlope} carries the paper's
whole lower-bound argument --- the $r=\kappa$ squeeze, the duality minimizing
slope, the positive-weight attainment and the completed square --- so none of
it is redone here.

\emph{Step B: the two one-sided slope comparisons.} The function $v$ satisfies
the hypotheses of \noderef{EnvelopeCompetitorBounds} with $I=\R$, and $W$
satisfies them with $I=(0,\infty)$ by \ref{it:cesl-Wcont} and
\ref{it:cesl-Wgrow}. The relevant times are $1+\ell>0$ for $\ell>0$ and
$1-\ell>0$ for $0<\ell<1$, so all four envelopes below are at admissible
parameters.

Let $\ell>0$. The competitor set defining $\Penv_\ell$ at $z_*$ in
\noderef{QuadEnvP} involves only values at the time $1+\ell>0$, where $W\le v$
by \ref{it:cesl-sand}; so for each $y\in F$,
\[
  W(y,1+\ell)-\frac{|y|^2}{2\ell}\ \le\ v(y,1+\ell)-\frac{|y|^2}{2\ell}
   \ \le\ \Penv_\ell v(z_*) ,
\]
the second inequality being the domination half of clause 1 of
\noderef{EnvelopeCompetitorBounds} for $v$. The extremality half of that clause
for $W$ then gives $\Penv_\ell W(z_*)\le\Penv_\ell v(z_*)$. Dividing by
$\ell>0$ after subtracting the values, which are both $0$ by
\ref{it:cesl-contact} and \eqref{eq:cesl-v0}, \noderef{EnvSlopeP} turns this
into
\begin{equation}\label{eq:cesl-plus}
  \Senv_\ell^+W(z_*)\le\Senv_\ell^+v(z_*)\qquad(\ell>0).
\end{equation}

Let $0<\ell<1$. The competitor set defining $\Menv_\ell$ at $z_*$ in
\noderef{QuadEnvM} involves only values at the time $1-\ell>0$, where again
$W\le v$; so for each $y\in F$,
\[
  \Menv_\ell W(z_*)\ \le\ W(y,1-\ell)+\frac{|y|^2}{2\ell}
   \ \le\ v(y,1-\ell)+\frac{|y|^2}{2\ell} ,
\]
the first inequality being the domination half of clause 2 of
\noderef{EnvelopeCompetitorBounds} for $W$. The extremality half of that clause
for $v$ then gives $\Menv_\ell W(z_*)\le\Menv_\ell v(z_*)$, and
\noderef{EnvSlopeM} with the two vanishing values turns this into
\begin{equation}\label{eq:cesl-minus}
  \Senv_\ell^-v(z_*)\le\Senv_\ell^-W(z_*)\qquad(0<\ell<1).
\end{equation}

\emph{Step C: $H$ against the slopes of $v$.} From \eqref{eq:cesl-ham} and
\eqref{eq:cesl-H}, $H$ is the infimum of the backward slopes, so
\begin{equation}\label{eq:cesl-lowerM}
  H\le\Senv_\ell^-v(z_*)\qquad(\ell>0).
\end{equation}
In the other direction, $H_v(z_*)$ is by \noderef{EnvHamiltonian} the infimum of
the \emph{forward} slopes of $v$ at $z_*$, so for every $\varepsilon>0$ there is
$\ell>0$ with
\begin{equation}\label{eq:cesl-infP}
  \Senv_\ell^+v(z_*)<H+\varepsilon .
\end{equation}
Indeed, were $H+\varepsilon\le\Senv_\ell^+v(z_*)$ for every $\ell>0$, then
$H+\varepsilon$ would be a lower bound for that family and hence
$H+\varepsilon\le H_v(z_*)=H$ by \eqref{eq:cesl-H}, contradicting
$\varepsilon>0$.

\emph{Step D: $H$ is a lower bound for the forward slopes of $W$.} By
\ref{it:cesl-Wcont}, \ref{it:cesl-Wgrow} and \ref{it:cesl-Wsub} the hypotheses
of \noderef{EnvelopeSlopeOrder} hold for $W$ on the open interval
$I=(0,\infty)$. Its first clause, whose hypothesis $1+\sigma>0$ for
$\sigma\in[0,\ell']$ is immediate, gives that
\begin{equation}\label{eq:cesl-mono}
  \Senv_\ell^+W(z_*)\le\Senv_{\ell'}^+W(z_*)\qquad(0<\ell\le\ell').
\end{equation}

Fix $h$ with $0<h<1$ and suppose, for contradiction, that
$\Senv_h^+W(z_*)<H$. Apply the second clause of \noderef{EnvelopeSlopeOrder} at
$z_*$ with this $h$ and with
$\varepsilon=(H-\Senv_h^+W(z_*))/2>0$; its hypothesis $1+\sigma>0$ for
$|\sigma|\le h$ holds because $h<1$. It supplies $\delta>0$ such that
$\Senv_\ell^-W(z_*)\le\Senv_h^+W(z_*)+\varepsilon$ whenever $0<\ell<\delta$.
Choose $\ell=\min\{\delta/2,1/2\}$, so $0<\ell<\delta$ and $\ell<1$. Then
\eqref{eq:cesl-lowerM}, \eqref{eq:cesl-minus} and that bound give
\[
  H\ \le\ \Senv_\ell^-v(z_*)\ \le\ \Senv_\ell^-W(z_*)
   \ \le\ \Senv_h^+W(z_*)+\frac{H-\Senv_h^+W(z_*)}{2}
   =\frac{H+\Senv_h^+W(z_*)}{2}\ <\ H ,
\]
the last inequality because $\Senv_h^+W(z_*)<H$. This contradiction shows
\begin{equation}\label{eq:cesl-low}
  H\le\Senv_h^+W(z_*)\qquad(0<h<1).
\end{equation}
For $h\ge1$ the same bound follows from \eqref{eq:cesl-low} at $1/2$ together
with the monotonicity \eqref{eq:cesl-mono} from $1/2$ to $h$. Hence
\begin{equation}\label{eq:cesl-lowAll}
  H\le\Senv_h^+W(z_*)\qquad(h>0).
\end{equation}

\emph{Step E: the limit.} We verify the two halves of the order
characterization of a limit for the function $\ell\mapsto\Senv_\ell^+W(z_*)$
along $\ell\downarrow0$.

If $a<H$, then by \eqref{eq:cesl-lowAll} $a<\Senv_\ell^+W(z_*)$ for
\emph{every} $\ell>0$, so this holds in particular for all $\ell$ near $0$ with
$\ell>0$.

If $a>H$, apply \eqref{eq:cesl-infP} with $\varepsilon=a-H>0$: there is
$\ell_0>0$ with $\Senv_{\ell_0}^+v(z_*)<H+(a-H)=a$, and then
\eqref{eq:cesl-plus} gives $\Senv_{\ell_0}^+W(z_*)<a$. For every $\ell$ with
$0<\ell<\ell_0$, the monotonicity \eqref{eq:cesl-mono} gives
$\Senv_\ell^+W(z_*)\le\Senv_{\ell_0}^+W(z_*)<a$; and the interval
$(0,\ell_0)$ is a neighbourhood of $0$ within $(0,\infty)$.

Therefore $\Senv_\ell^+W(z_*)\to H$ as $\ell\downarrow0$, which with
\eqref{eq:cesl-H} is the second conclusion, for the same $H$.
\end{proof}
